\documentclass[11pt,a4paper]{article}
\usepackage{amsmath,amssymb,amsthm,bm}
\usepackage{geometry}
\geometry{margin=1in}
\usepackage{hyperref}
\usepackage{graphicx}
\usepackage{setspace}
\usepackage{booktabs}
\onehalfspacing

\title{\textbf{Paper OOL-2: Iron--Sulfur Redox and Mineral Constraint Scaffolds:\\ The First Geochemical Constraint Layer}}
\author{Steve Bico Mujjabi, MD\\
\small Independent Researcher\\
\small ORCID: \href{https://orcid.org/0009-0001-0556-5516}{0009-0001-0556-5516}}
\date{May 2026}

\begin{document}
\maketitle

\begin{abstract}
This paper begins the concrete Part II chain with iron--sulfur redox and mineral constraint scaffolds. Fe--S surfaces are treated as localized energy-processing boundaries: they can host electron transfer, surface catalysis, pore-scale confinement, and continuity with later Fe--S cluster biochemistry. The claim is conditional, not triumphalist: mineral Fe--S matters for origins only when a planet supplies sustained donors, acceptors, renewal, and retention. The companion script illustrates boundary-local flux amplification in a toy CDFD field and is explicitly not a molecular simulation.
\end{abstract}

\section{Introduction}
The origin of life is inseparable from the problem of maintaining chemical work against dilution and side reactions. Redox chemistry provides the elementary mechanism: electrons move between donors and acceptors, coupling environmental gradients to bond rearrangement. Modern biochemistry places ancient fingerprints on this problem: iron--sulfur clusters remain central to electron transfer in respiration and photosynthesis, motivating the hypothesis that mineral Fe--S phases participated in the earliest energy-processing networks~\cite{Russell1994,Martin2007,Branscomb2018}.

Paper OOL-2 opens the concrete geochemical sequence by asking a precise question: \textbf{can Fe--S \emph{geometry} and connectivity be represented as organizers of redox throughput, not merely as passive catalysts?} We connect this question to hydrothermal petrology and to a phenomenological transport field in which gradients in constraint modulate flow. Paper~OOL-3 extends the same viewpoint to distributed conduction (magnetite-like networks).

\subsection*{Bridging gap from prior work}
The present series treats prebiotic stages as a ladder of transport and constraint mechanisms. No prior installment exists; the unresolved gap is empirical: which mineralogies and fluid chemistries first produced \emph{persistent} redox channeling under open-system forcing.

\textbf{Status: Inferred}

\paragraph{Capstone channel mapping.}
Paper 12 treats Fe--S redox gradients as one early implementation of the energy-input function $I_E$. Later papers allow other implementations, including photochemical materials, but this first stage does not assume chlorophyll or any mature photosynthetic machinery. The purpose here is narrower: establish how mineral geometry can localize usable redox work.

\section{Model and Assumptions}

\subsection{Notation (CDFD alignment)}
Throughout the OOL series, $C$ denotes a nonnegative \emph{constraint} or transport-resistance field (earlier drafts used $\Lambda$ for the same role). $\Phi$ is an intensity variable for redox driving (electrochemical potential or associated flux potential). The dimensionless ratio
\begin{equation}
\Psi_s \;=\; \frac{\Phi}{C}
\end{equation}
measures local ``overload'' of flow relative to resistance. Species concentrations in reaction kinetics are written $c_i$ when needed, to avoid clashing with $C$.

\textbf{Status: Derived}

\subsection{Geochemical formation of Fe--S assemblages}
Hydrothermal fluids rich in Fe$^{2+}$ and reduced sulfur precipitate Fe--S phases; representative steps include
\begin{align}
\text{Fe}^{2+} + \text{H}_2\text{S} &\rightarrow \text{FeS} + 2\text{H}^+,\\
\text{FeS} + \text{S} &\rightarrow \text{FeS}_2.
\end{align}
Mackinawite, pyrite, and related assemblages often grow as porous networks rather than isolated grains, yielding high surface area, tortuous transport, and internal confinement~\cite{Martin2007,Hazen2008}. These are the geological preconditions for localized chemistry.

\textbf{Status: Inferred}

\subsection{Redox thermodynamics}
Iron cycles between oxidation states,
\begin{equation}
\text{Fe}^{2+} \leftrightarrow \text{Fe}^{3+} + e^-,
\end{equation}
with electrochemical work set by Nernst-type relations; a compact statement is $\Delta G = -n F \Delta E$ for $n$ electrons transferred, Faraday $F$, and potential gap $\Delta E$. Fe--S minerals scaffold this chemistry at interfaces where water, dissolved gases, and mineral terminations coexist.

\textbf{Status: Derived}

\subsection{Electron transport as a field process}
Redox events embed in a continuum description of electron transport. We write the current density in the resistance form
\begin{equation}
\bm{J}_e = -\frac{1}{C}\nabla\Phi,
\label{eq:Je}
\end{equation}
so that large $C$ (high resistance) suppresses flow at fixed potential gradient. In Fe--S environments, $\nabla\Phi$ arises from vent chemistry and fluid mixing; $C$ encodes mineral connectivity, disorder, and interfacial barriers. Eq.~\eqref{eq:Je} is \emph{phenomenological}: it does not resolve hopping microphysics or Marcus rates.

Rather than viewing redox reactions as isolated events, we interpret them as part of a continuous transport process. This formulation highlights that electron flow is not only driven by gradients but also fundamentally shaped by the structure of the medium.

\textbf{Status: Hypothesis}

\subsection{Catalysis and confinement}
Surfaces lower activation barriers via adsorption and transition-state stabilization; schematically $A+B \rightarrow AB$ is accelerated at the interface. These surfaces adsorb reactants, stabilize intermediate states, and increase effective collision rates, enabling reactions that would otherwise be kinetically unfavorable in bulk solution.

Porous Fe--S masses retard diffusion, increasing residence time. A Fickian picture with spatially varying resistance is
\begin{equation}
\frac{\partial \rho}{\partial t} = \nabla\cdot\left(\frac{1}{C}\nabla\rho\right) + R,
\end{equation}
where $\rho$ is a scalar concentration proxy and $R$ gathers reactions. In regions where $C$ is large (such as pore throats), diffusion is reduced, reactants remain localized, and reaction probability increases. This prevents the rapid dilution of chemical species, a critical requirement for sustained chemical networks.

\textbf{Status: Inferred}

\section{Main Results}

\subsection{Conceptual synthesis}
Within the assumptions above, Fe--S systems couple three roles: (i)~they generate and maintain redox gradients; (ii)~they channel electron transport through structured resistance; (iii)~they confine solutes so that surface catalysis can outperform bulk washout.

The key insight is that Fe--S systems do not merely catalyze reactions; they organize energy flow. Electron transport and catalytic activity are deeply coupled: transport distributes energy via $J_e = -\frac{1}{C} \nabla \Phi$, while catalysis transforms matter. Together, these features motivate treating vent-associated Fe--S as the first geological ``energy-processing substrate'' in the series narrative, prior to extended magnetite networks (Paper 3) and interfacial water (Paper 4).

\textbf{Status: Inferred}

\subsection{Continuity with biochemistry (interpretive)}
Biological Fe--S clusters perform electron transfer in contemporary enzymes; mineral Fe--S offers a geochemical precedent for localized, mixed-valence chemistry~\cite{Russell1994}. This continuity argument is qualitative: homology of function does not establish homology of historical pathway.

\textbf{Status: Hypothesis}

\section{Numerical Verification}
Companion script: \texttt{supplementary\_ool\_02.py}. The implementation places a high-$C$ mineral block in a 2D grid, applies a simple Laplacian-like update to a nonnegative ``redox flux'' field $\Phi$, and boosts $\Phi$ where $|\nabla C|$ exceeds a threshold (interface proxy for catalytic acceleration). This is a \emph{toy} illustration: it is not a molecular simulation of mackinawite surface chemistry.

Observed outputs (fixed seed in script):
\begin{itemize}
\item Initial mean redox flow: $0.149$
\item Final mean redox flow: $50.224$
\item Peak interface flow: $100.000$ (the imposed numerical cap)
\end{itemize}

The mean and peak amplification support the claim that, \emph{in this model}, mineral boundaries with steep constraint gradients can sustain elevated local throughput relative to bulk. The peak value reaches the script cap, so the result should be read as a saturation warning as well as an amplification signal.

\textbf{Status: Derived}

\begin{figure}[ht]
\centering
\includegraphics[width=\textwidth]{outputs/paper_02/fe_s_redox_interface.png}
\caption{Toy Fe--S interface diagnostic. The left panel shows the final redox-flux proxy with the mineral boundary overlaid; the right panel shows mean flux and peak interface flux across the run. The figure is an illustration of localization under a constraint gradient, not a molecular simulation.}
\end{figure}

\section{Geochemical Boundary-Condition Ledger}
Fe--S chemistry is not a free-floating catalyst list. It becomes origin-of-life relevant only when embedded in a fed environment. In the present paper the required boundary conditions are:
\begin{itemize}
\item a sustained donor side, such as Fe$^{2+}$, H$_2$, H$_2$S, or reduced vent fluids;
\item an acceptor side, such as CO$_2$, oxidized mineral phases, or other environmental oxidants;
\item a pH or redox discontinuity that keeps $\nabla\Phi$ from immediately relaxing;
\item porous retention so intermediates survive longer than the reaction time;
\item repeated replenishment, precipitation, or fracture flow so the mineral does not become a closed dead surface.
\end{itemize}
This list prevents the paper from claiming that Fe--S minerals alone make life. They are a scaffold only when the surrounding planet supplies drive, disequilibrium, and renewal.

\textbf{Status: Inferred}

\section{Failure Modes}
The Fe--S step fails if any of three bottlenecks dominates. First, the system may have surface catalysis without a persistent gradient; this gives transient chemistry but no maintained flux. Second, the gradient may exist but products may be lost faster than they react; this motivates the compartment paper. Third, the mineral may bind molecules too strongly or poison its own sites, converting selection into immobilization. These failures are not weaknesses of the series; they are the reason the next papers are necessary.

\textbf{Status: Derived}

\section{Experimental and Environmental Anchors}
The Fe--S claim should be attached to real experimental and geological anchors rather than treated as a symbolic first step. Alkaline hydrothermal-vent models provide a natural reason to combine redox gradients, mineral precipitation, pH contrast, and pore-scale confinement~\cite{Russell1994,Martin2007,Sojo2016}. A recent redox-focused review also supports treating early Earth as a multi-environment redox system rather than a single-location story~\cite{Moldogazieva2026}. Laboratory vent-reactor work attempts to reproduce some of those coupled conditions by growing mineral membranes under controlled disequilibria~\cite{Barge2014}. These experiments are not a proof of life emergence, but they show that the relevant variables can be made concrete: fluid composition, pH, membrane mineralogy, redox potential, pore growth, and transport.

The framework also needs to acknowledge counterarguments. Jackson argued that natural pH gradients in alkaline vents may not have been sufficient or relevant in the way some origin scenarios require~\cite{Jackson2016}. This objection is valuable for CDFD because it prevents the theory from hiding behind a single favored environment. The model should therefore be tested across a family of fed mineral reactors, not only one idealized vent geometry.

\textbf{Status: Inferred}

\section{Fe--S Clusters as a Later Continuity Signal}
Modern Fe--S clusters are not evidence that free mineral surfaces were the first catalysts. They are, however, a continuity signal: life still uses iron, sulfur, redox hopping, and ligand-stabilized electron transfer in core metabolism. Prebiotic synthesis of Fe--S-like clusters under photochemical conditions further widens the possible bridge between geochemistry and cofactors~\cite{Bonfio2017}. The conservative CDFD claim is therefore:
\begin{equation}
\text{mineral Fe--S function} \rightarrow
\text{ligand-bound Fe--S function}
\rightarrow
\text{enzyme-embedded Fe--S function}.
\end{equation}
The arrow is functional, not genealogical. It says the same physical problem---localized, reversible electron transfer---appears at multiple later levels of biological organization.

\textbf{Status: Inferred}

\section{Interpretation}
The numerical pattern is consistent with the qualitative geochemical picture: interfaces organize flow. The model does not, however, establish which specific Archean mineral facies dominated, nor does it compute absolute reaction rates for prebiotic syntheses~\cite{Huber1997,Wachtershauser1988}. Interpretations of ``first metabolism'' therefore remain at the Hypothesis/Inferred tier pending tighter coupling to kinetics and fluid dynamics.

\textbf{Status: Inferred}

\section{Limitations and Open Problems}
\begin{itemize}
\item Interface catalysis is injected via a gradient threshold rule rather than ab initio activation energies.
\item Polymorph-specific Fe--S properties (mackinawite vs pyrite electronic structure) are not distinguished.
\item Natural porous networks require topology statistics beyond a single square inclusion.
\item Three-dimensional fracture--vein systems and coupled fluid flow are not modeled here.
\end{itemize}

\textbf{Status: Open}

\section{Conclusion}
Fe--S-rich hydrothermal environments provide a geochemically motivated stage 1 for the OOL series: redox energy, catalytic surfaces, and confinement map naturally onto a constraint-field description with $\Psi_s=(\Phi/C) \cdot S \cdot M_s$. Companion numerics illustrate boundary-local amplification in a controlled idealization; Paper 3 extends the transport story to distributed, mixed-valence mineral conduits.

\textbf{Status: Inferred}

\section{Reproducibility Note}
\begin{verbatim}
python supplementary_ool_02.py
\end{verbatim}
Uses NumPy only; runtime sub-second on typical laptops. The script writes compact reproducibility outputs under \texttt{outputs/paper\_02/}.
\bibliographystyle{plain}
\bibliography{ool_references}
\end{document}
